\section{From local plane energies to the sphere}\label{sec:sphere}

The energy identification of Section~\ref{sec:energy} was obtained in
ordinary patches of a quantum cone.  We now transfer it to the fixed-area
sphere and include the point at which the contour begins and ends.  The
output is a global variational statement in the same coupling as the
metric-measure convergence.  This common coupling will later allow us to
compare the walks and their metric spaces simultaneously.

Write $\dd A$ for round area on the conformal sphere, and set
\[
 \mathcal D(u)=\int_{S^2}|\nabla u|^2\,\dd A
 \quad\text{for }u\in C(S^2)\cap W^{1,2}(S^2),
\]
with $\mathcal D(u)=+\infty$ for other continuous functions.  The same
notation on an open coordinate domain means the integral over that
domain, with local weak differentiability required.  Conformal
invariance makes these definitions independent of the chosen chart.

\subsection{Transfer on the interior of a contour interval}

First extend the plane conclusion beyond a single ordinary patch.
Retain all forward dyadic unit changes $v=2^k$ from
Proposition~\ref{inp:transport}.  For each fixed $k$, the transformed
rate still tends to infinity, so Theorem~\ref{en:variational} applies
on the ordinary patches inside that cone's normalization disk.
Under the joint zoom \eqref{inp:joint-zoom}, these disks exhaust the
plane away from the cone point.  Pulling the countably many
probability-one conclusions back to the original coordinate proves
the local energy identity everywhere off that point.  Both the
network costs and the two-dimensional Dirichlet integral are unchanged
by these transformations.

For $0<d<1/2$, let $I_d=[d,1-d]$.  The decorated interval surface
associated with the contour on $I_d$ has an interior determined by those
increments.  A compact subset of that interior, together with a small
collar, has all its contour visits strictly inside $I_d$.  Thus its
eventual discrete incidences are read from the internal word; no
information from either omitted end of the tour is needed.  These are
the interval-surface and complete-core assertions used in
Lemma~\ref{inp:complete-core}.

Here is the part of the excursion comparison that permits new
variational observations.  It concerns the joint law of the arrays
and the word, not independence after weak convergence.

\begin{lemma}[Additional observations under excursion conditioning]
\label{sph:block-density}
Fix $0<a<b<1$.  Let $W_N$ be the diffusively rescaled free axial
increment path on $[\lfloor aN\rfloor,\lfloor bN\rfloor]$, with
linear interpolation, and let $Z_N$ be any measurable function of
these increments with values in a fixed compact metrizable space.
Suppose $(W_N,Z_N)\Rightarrow(W,Z)$ along a sequence of even $N$.
There is a nonnegative continuous function $f_{a,b}$ of the limiting
path, with expectation one, such that the corresponding pair under
quadrant-excursion conditioning converges to the law
\[
 f_{a,b}(w)\,\Law(W,Z)(\dd w,\dd z).
\]
In particular, every probability-one property of $(W,Z)$ also
holds for this conditioned limit.  A countable product of
compact-valued observations is allowed for $Z_N$.
\end{lemma}

\begin{proof}
The lemma \emph{Excursion densities} in the spanning-tree sphere
section of \cite{MetricMaps} proves the following statements for
the density $f_N$ of this increment block.  If $w_N\to w$
uniformly, then $f_N(w_N)\to f_{a,b}(w)$; the limit depends
continuously on the two coordinate minima and the net increment.
Moreover, $\sup_N\E[f_N(W_N)^2]<\infty$ under the free law.
These assertions follow there from the entrance-and-exit killed-walk
density and its Gaussian bounds.

Realize the assumed joint convergence almost surely.  For any bounded
continuous test $\Phi$, first truncate the densities at height $K$.
Sequential convergence and bounded convergence give convergence of
the weighted expectations with truncated densities.  The error in
removing truncation is at most
$K^{-1}\|\Phi\|_\infty\sup_N\E[f_N(W_N)^2]$, and the analogous
limit error has the same bound.  Hence
\[
 \E[f_N(W_N)\Phi(W_N,Z_N)]
 \longrightarrow\E[f_{a,b}(W)\Phi(W,Z)].
\]
Taking $\Phi=1$ gives total mass one.  The finite-resolution
change of density depends only on the increment block, so its
weighted expectation is precisely that of the additional block
data under excursion conditioning.  This proves the assertion.
\end{proof}

\begin{proposition}[Energy identification under the sphere conditioning]
\label{sph:transfer}
Consider any joint subsequential extraction of the sphere contour and
the countable internal variational arrays on $I_d$, for every rational
$d\in(0,1/2)$, using the extraction convention of
Proposition~\ref{inp:array-interface}.  Almost surely, on every open
domain compactly contained in the interior of one of these interval
surfaces, both primal and dual inner relaxed energies equal
$\mathcal D$.  This identity holds simultaneously for all continuous
targets, with the value $+\infty$ when the target is not locally
Sobolev or its Dirichlet integral diverges.
\end{proposition}

\begin{proof}
At each resolution retain the energy minima with rational value
constraints and finite unions of relative-time intervals as selectors.
Only increments within $I_d$ enter these internal graphs and tests.
Their compactified values are therefore permissible additional block
data in Lemma~\ref{sph:block-density}.  The corresponding free
experiment is the bilateral sewing at rate $m=2n$.  In that experiment
adjoin these jump-indexed internal tests to the Poisson-instruction
arrays before extraction.  Each added test is computed from its exact
integer increment block and is unchanged when an independent Poisson
clock is attached to the word.  The earlier local identity applies
in this enlarged extraction.  To read it from the internal tests,
use complete pieces and fixed spatial collars as described below.
The vanishing discrepancy of the rescaled clocks only
locates the selectors inside these collars; the numerical energy
outputs themselves are retained unchanged, without a continuity
claim for them under relabeling.

Starting with the given conditioned extraction, take a further
subsequence on which the corresponding free block and its internal
array converge jointly.  This is possible by tightness of the free
contour and compactness of the array space, simultaneously for the
countably many intervals by diagonal extraction.  The weighted-limit
lemma then identifies the already specified conditioned block limit
with the density tilt of this free joint limit.

Theorem~\ref{en:variational} identifies the inner relaxed energies in
that bilateral experiment, away from its distinguished point.  The
distinguished point does not belong to this interval surface:
time zero is almost surely a singleton Brownian contour fiber.
Indeed, each coordinate takes values strictly below its value at zero
arbitrarily close to zero on both sides.  Thus no nontrivial
identification link in either coordinate can have zero as an
endpoint, and neither can the equivalence relation generated by
those links identify zero with another time.  Since $I_d\subset(0,1)$,
its image excludes the distinguished point.

To express the local energy conclusion in terms of the block alone,
reconstruct its abstract decorated interval surface $H_d$ from the
limiting increments using the interval reconstruction in
Proposition~\ref{inp:sewing}.  This reconstruction supplies the surface
and its restricted curve modulo conformal coordinates.  Evaluate
the local identity in a conformal representative, using a countable
atlas and compact exhaustions of its interior; conformal invariance
makes the event independent of that representative.  Both the
restricted curve locating the labels and the coordinate domains for
the targets belong to this reconstructed interval surface.
For a compact $K$ in one such domain and an open enlargement
$U$ with a collar, Lemma~\ref{inp:complete-core} supplies all incidences
near $K$ from the internal word.  The selectors of
Lemma~\ref{inp:sandwich} put every constraint on $K$ between
constraints on internal graphs contained in $U$.  Rational tolerances
and compact exhaustion therefore express the inner relaxed energy
using only the limiting block and its retained arrays.  A change of
the conformal representative preserves both the transported energy
identity and the Dirichlet integral.  The resulting assertion is
intrinsic to $H_d$; the surrounding cone or sphere embedding supplies
no further data for its definition.

The simultaneous quantifier on targets can also be retained in this
transfer.  On a compact coordinate set use a countable dense library
of continuous functions and rational tolerances.  Uniform closeness
of two targets enlarges or shrinks a tolerance constraint by at most
that distance.  The library therefore describes all the relaxed
value tests, jointly measurably in the retained arrays and the target.
This is a discretization of constraints, not an approximation of
Dirichlet energy merely from values on a uniformly dense library.
For varying charts use a countable atlas and compact exhaustions.
The usual Dirichlet integral, with its infinite-value convention,
is likewise Borel on continuous targets: weak derivatives and their
energy can be tested against countably many smooth functions on the
exhaustion.  Existence of a continuous target violating the identity
is consequently an analytic event, obtained by projection from its
Polish function space.  The event of equality for all targets is
universally measurable.  Theorem~\ref{en:variational} gives precisely
this simultaneous event under the free law, so absolute continuity
applies to it without a target-dependent exceptional set.

Absolute continuity of the joint block law now transfers the
probability-one identity to the excursion experiment.  A countable
intersection over $d$, charts, and exhaustions proves the result.
The rest of the sphere contour may be retained in this coupling: the
complement of the event already has zero probability in the internal
marginal, hence also in the joint law.
\end{proof}

The interval interiors exhaust the sphere away from the tour endpoint,
denoted by $p$.  In fact, the inverse image under the continuous contour
of a compact set disjoint from $p$ is compact and avoids times $0$ and
$1$; the same holds after a sufficiently small spatial enlargement.
The interval-surface identification then places the set and its collar
in an interior with $d>0$.  Consequently
Proposition~\ref{sph:transfer} gives the local energy identity at every
point other than $p$.

\subsection{Removing the tour endpoint}

There are two different issues at $p$.  A limit of bounded-energy
functions must have no extra Sobolev singularity there, and recovery
sequences must incur no extra discrete energy there.  The following
elementary fact handles both.

\begin{lemma}[Continuous Sobolev removability]\label{sph:removability}
Let $p$ be a point of a smooth two-dimensional surface.
If $u$ is continuous near $p$, is locally $W^{1,2}$ off $p$, and has
finite Dirichlet integral near $p$, then $u$ is $W^{1,2}$ across $p$.
Moreover, if $u$ is continuous and Sobolev on a compact surface, there
are continuous Sobolev functions $u_r$, constant on a neighborhood of
$p$ and equal to $u$ outside a shrinking disk, such that
\[
 \|u_r-u\|_\infty\longrightarrow0,
 \qquad
 \|u_r-u\|_{W^{1,2}}\longrightarrow0.
\]
\end{lemma}

\begin{proof}
Work in a coordinate disk centered at $p$.  A point has zero
two-dimensional Sobolev capacity: radial logarithmic cutoffs equal to
one on a shrinking inner disk and zero outside a fixed disk have
Dirichlet energies tending to zero.  Testing the weak-derivative
identity off $p$ with the complementary cutoffs proves the identity
across $p$.  More explicitly, the additional term containing the
cutoff gradient tends to zero by Cauchy--Schwarz, since $u$ is bounded
on the support of a fixed smooth test function, while the term
containing $\nabla u$ converges by its local square integrability.

For the approximation choose a smooth cutoff $\chi_r$ equal to one
on $B(p,r)$, zero outside $B(p,2r)$, taking values in $[0,1]$, and
with $\int|\nabla\chi_r|^2\,\dd A\le C$ independently of $r$.
Set
\[
 u_r=u(p)+(1-\chi_r)(u-u(p)).
\]
The uniform error is bounded by the oscillation of $u$ on $B(p,2r)$.
For the gradient error,
\[
 \int|\nabla(u_r-u)|^2\,\dd A
 \le 2\int_{B(p,2r)}|\nabla u|^2\,\dd A
       +2C\sup_{B(p,2r)}|u-u(p)|^2,
\]
and both terms tend to zero.  The ordinary $L^2$ error follows from
uniform convergence.
\end{proof}

In a coupled extraction, choose a location $\pi_n(v)\in S^2$ for each
primal vertex from one of its contour visits; different choices have
uniformly vanishing discrepancy.  Use $\pi_n^*$ on the dual.
For functions on the varying vertex sets, uniform convergence to $u$
means
\[
 \max_{v\in V_n}|u_n(v)-u(\pi_n(v))|\longrightarrow0.
\]
The small-image joining and mesh assertions in the input guarantee that
this convention agrees with the local comparisons already used.
Harmonicity on an open set $U$ means that the graph Laplacian vanishes
at every vertex whose location is in $U$.  For a compact subset of
$U$, all adjoining stars also lie in $U$ once $n$ is large; thus
this convention permits the buffered uses of
Lemma~\ref{en:harmonic-compactness}.

\begin{theorem}[Global sphere energies]\label{sph:energy}
From any subsequence of the tree-decorated sphere maps, one can take a
further joint extraction and a coupling with the limiting ordinary
quantum sphere such that the following statements hold almost surely.
\begin{enumerate}
\item If $u_n$ converges uniformly to $u\in C(S^2)$, then
\[
 \liminf_{n\to\infty}\en_n(u_n)\ge\mathcal D(u).
\]
\item For each $u\in C(S^2)\cap W^{1,2}(S^2)$ there are uniformly
convergent approximations $u_n$ for which
\[
 \en_n(u_n)\longrightarrow\mathcal D(u).
\]
Both assertions hold on the dual with $\en_n^*$ and $\pi_n^*$.
\item On every open set, graph-harmonic functions with bounded energy
and bounded absolute value are locally equicontinuous under these
locations and have locally uniformly convergent subsequences.  This
holds also on the dual and at the tour endpoint.
\end{enumerate}
All these conclusions may be required in the same extraction as the
metric and measure comparison below.
\end{theorem}

\begin{proof}
Retain the internal tests of Proposition~\ref{sph:transfer} and also
the whole-sphere tests defining uniform relaxed energy.  The latter
are countable compactified arrays; they need no local conditional
kernel.  The selector sandwich applies to them in the same way.

Suppose first that $\liminf_n\en_n(u_n)<\infty$ and pass to a subsequence
realizing this liminf.  On disjoint compact-interior patches away from
$p$, Proposition~\ref{sph:transfer} bounds their limiting Dirichlet
integrals by the discrete energies in those patches.  Exhausting the
complement of $p$ gives
\[
 \int_{S^2\setminus\{p\}}|\nabla u|^2\,\dd A
 \le\liminf_n\en_n(u_n).
\]
Lemma~\ref{sph:removability} supplies the weak derivative at $p$.
This proves the first assertion, including its infinite-value case.

For the upper bound, first take a continuous Sobolev target that is
constant near $p$.  Cover the complement of a smaller constant disk
by finitely many patches from Proposition~\ref{sph:transfer}.
Subdivide it into finitely many smaller coordinate regions with
piecewise smooth boundaries, each contained in one patch, and thicken
these regions slightly.  The boundaries have zero round area, so
their collars can be chosen with arbitrarily small total integral
of $|\nabla u|^2$.  Local recovery on the thickened regions and
successive applications of Lemma~\ref{en:fundamental} now give
recovery on the compact complement: the main energy is bounded by
the sum of the target energies on those regions, and only the
chosen small collars are counted more than once.  The error terms
from gluing tend to zero by uniform agreement on each fixed
collar.  This proves recovery with arbitrarily small energy slack.
Extend by the constant on the
smaller disk.  To make the extension exact, use a cutoff whose
transition lies in a fixed annulus where the target is constant.
Apply the finite-cover recovery just proved to a slightly larger
annular neighborhood and a smooth cutoff $b$ that is constant on
its two collars.  If $b_n$ is its recovery and
$\delta_n>\|b_n-b\|_\infty$ tends to zero, replace $b_n$ by
\[
 \left(\frac{b_n-\delta_n}{1-2\delta_n}\right)\vee0\wedge1.
\]
This is exactly zero and one on the respective flat collars, and
its energy is at most $(1-2\delta_n)^{-2}\en_n(b_n)$.
The functions being joined differ by $o(1)$
on the transition, so the extra term in the gluing estimate is $o(1)$.
The mesh tends to zero, hence every edge crossing either gluing
boundary is contained in the fixed collar for all sufficiently large
$n$.

For a general continuous Sobolev $u$, apply the approximation
$u_r$ of Lemma~\ref{sph:removability}.  Recover each fixed $u_r$ first,
then use a diagonal choice of $r\downarrow0$.  Uniform convergence and
convergence of $\mathcal D(u_r)$ give the desired upper bound; the
first assertion makes it an equality.  The proof for the dual is
identical.

Away from $p$, harmonic compactness follows from
Lemma~\ref{en:harmonic-compactness}.  To include $p$, take any fixed
finite stack of disjoint annuli centered at $p$, each of the same
aspect ratio, in a domain of harmonicity.  Every transition annulus
avoids $p$.  By the preceding local recovery, its dual cutoff can
be taken with energy bounded by a constant depending only on the
aspect ratio, with exact constant values on its two sides.  Fill
the inner side by the constant; this uses only a flat collar in the
annulus.  Lemma~\ref{en:circuit} gives a primal circuit in each annulus
whose oscillation squared is bounded by that constant times the
primal energy in the annulus.  Among $k$ disjoint annuli, one
therefore has oscillation at most $C k^{-1/2}$ when the total energy
is bounded.  The maximum principle transfers this oscillation to
all vertices inside its circuit.  First choose $k$, and then take
large $n$; no annulus radius depends on $n$.  Finite covers give
equicontinuity, and bounded values give compactness.  Exchanging
primal and dual proves the remaining assertion.
\end{proof}

\subsection{Keeping the metric and the quantum area in the same coupling}

The coordinate used for energy convergence must also describe the
metric-measure limit.  The contour correspondence makes this precise.

\begin{proposition}[Joint spatial comparison]\label{sph:comparison}
The extraction in Theorem~\ref{sph:energy} can be chosen so that its
limiting sphere is the one in the metric-measure limit and, almost
surely,
\begin{align}
 \sup_{v,w\in V_n}
 \big|a_n d_n(v,w)-D_h(\pi_n(v),\pi_n(w))\big|&\longrightarrow0,
 \label{sph:distortion}\\
 \sup_{x\in S^2}\dist_{D_h}(x,\pi_n(V_n))&\longrightarrow0,
 \label{sph:dense}\\
 (\pi_n)_*\mu_n&\Longrightarrow\mu_h.
 \label{sph:measure}
\end{align}
The maximum round diameter of a projected edge tends to zero.
The measure $\mu_h$ is nonatomic and has full support, and $D_h$
induces the round topology.
\end{proposition}

\begin{proof}
The sphere-identification proposition and the final normalization
argument for spanning-tree maps in \cite{MetricMaps} give, jointly
with the excursion, the contour correspondence
\[
 \bigl\{(v_{\lfloor2nt\rfloor},\eta(t)):0\le t\le1\bigr\}
\]
with distortion tending to zero.  Here $v_j$ is the primal vertex
at the start of tour step $j$ and $\eta$ is the sphere contour.
Those results identify the metric of this very peanosphere limit;
they do not introduce an independent quantum sphere.  Their
deterministic diameter normalization is the $a_n$ in
Theorem~\ref{thm:main}.  Choosing one visit for each vertex changes
the correspondence by a uniformly vanishing discrepancy, which
gives \eqref{sph:distortion} and \eqref{sph:dense}.

There are $2n$ tour step starts.  Their vertex occupation measure is
$\mu_n$: horizontal tree traversals are counted at their starting
vertices, and vertical steps count the non-tree half-edges at their
incident vertices.  Every half-edge is counted once, including both
half-edges of a loop.  The area parameterization of $\eta$ therefore
gives \eqref{sph:measure}.  Mesh convergence is the sphere
projection construction preceding Proposition~\ref{inp:topology}.
Nonatomicity, full support, and agreement of the metric and conformal
topologies are the quantum-area and metric properties in the same
repository input.

All these variables may be retained together with the compactified
energy arrays before extraction.  The metric identification is
unchanged upon a further extraction, so it holds in the coupling of
Theorem~\ref{sph:energy}.
\end{proof}
